\section{Trash}
\label{sec:trash}

Methods based on equivalent layer for dealing with total-field anomaly data 
have been traditionally developed in space domain 
\cite[e.g.,][]{silva_1986, oliveirajr-etal_2013}
by numerically solving the integral equation \ref{eq:B-t-generic-layer} 
for the unknown physical-property distribution $p(\mathbf{r}'')$, 
with predefined constant values for the inclination $I = I_{p}$ and declination $D = D_{p}$. 
Those methods usually assume that total-field anomaly data 
can be satisfactorily approximated by the harmonic function $B_{t}(\mathbf{r})$
(equation \ref{eq:B-t-generic-layer}).
After obtaining a numerical solution for the physical-property distribution 
$p(\mathbf{r}'')$, it is used to compute a potential-field transformation. 
If the desired transformation does not
cause phase changes in the magnetic data, such as the upward/downward continuation, 
then the predefined constant values for the inclination $I_{p}$ and declination $D_{p}$ 
may be arbitrarily chosen. 
Otherwise, the $I_{p}$ and $D_{p}$ must assume the same constant inclination and declination
values defining the total-magnetization direction of the true sources. An example of 
transformation that produces phase changes in the magnetic data is the RTP.